\begin{textsample}{POS Dim 1 – ma – Score 42.00 – ma/ma\_em\_48.txt}  \label{ex:f1_pos_067}
Itinerário de Ciências Sociais, Econômicas e Administrativas Enfoque do itinerário : Filosofia | Geografia | História | Sociologia | Matemática As constantes mudanças trazem incertezas quanto aos aspectos econômicos, políticos e sociais \textbf{que} se misturam a \textbf{um} cotidiano \textbf{que} precisa ser captado por \textbf{um} sujeito cada vez mais transdisciplinar e ciente do seu papel \textbf{enquanto} agente transformador.
É certo \textbf{que} praticamente tudo \textbf{que} está à nossa volta e \textbf{que} utilizamos na escola, em casa ou no trabalho foi construído pelo \textbf{homem}, e esse \textbf{sujeito}, seja por ação individual ou coletiva, é capaz de transformar a sociedade conforme a época \textbf{que} esteja vivenciando.
Nesse aspecto, o itinerário de ciências sociais, econômicas e administrativas se apresenta como \textbf{uma} opção para o jovem \textbf{que} pretende seguir os cursos \textbf{que} englobam \textbf{uma} série de possibilidades de estudos perceptíveis e analisáveis à sua volta, como os mais diversos fenômenos sociais \textbf{que} se desenvolvem no seu contexto.
Podemos exemplificar assuntos sociais, econômicos e administrativos relevantes \textbf{que} podem percorrer os eixos estruturantes propostos para o itinerário.
Assuntos como gestão pública e privada, desenvolvimento e subdesenvolvimento, taxas de juros, aplicações financeiras e sua rentabilidade e liquidez, impostos, consumo, desemprego, inflação, fome, desenvolvimento sustentável e comércio internacional aprofundam aprendizagens importantes nas áreas de Ciências Humanas e Sociais Aplicadas e de Matemática, retomando e aprofundando temas abordados na BNCC, com o intuito de situar esses \textbf{sujeitos} no tempo e espaço em \textbf{que} estão inseridos, orientando -os à continuidade de reflexão e abstração com perspectivas conscientes de mediação e intervenção social, de investigação científica, de empreendedorismo, além de infinitas possibilidades criativas.
Vale observar \textbf{que} o pleno desenvolvimento da pessoa, seu preparo para o exercício da cidadania e sua qualificação para o trabalho são prerrogativas do ensino médio, o \textbf{que} reforça a contribuição deste itinerário no sentido de garantir subsídios na preparação desse jovem protagonista em meio aos seus dilemas enfrentados diariamente.
A ciência econômica lida com a disposição de recursos limitados \textbf{que} obrigam tanto grandes estadistas como pequenos consumidores a estarem diariamente optando como serão gastos seus recursos.
Entender o cenário econômico no qual está inserido vai da análise comportamental do indivíduo até grandes decisões coletivas, \textbf{que} envolvem toda \textbf{uma} sociedade.
Assim, identificar e analisar o contexto histórico no qual estamos inseridos também contribui para a formação cidadã.
Da mesma forma, administração amplia e possibilita o entendimento em negócios, pessoas ou recursos, com o objetivo de alcançar metas definidas, e, para isso, a matemática garante elementos sustentadores na investigação científica dos fenômenos.
No \textbf{que} tange aos aspectos específicos para a \textbf{abordagem} do itinerário de Ciências Sociais, Econômicas e Administrativas, o enfoque está fundamentado em aprendizagens dos componentes curriculares filosofia, geografia, história, sociologia e matemática, aprofundando competências e habilidades relacionadas ao itinerário articuladas aos temas \textbf{contemporâneos}.
O presente itinerário \textbf{fomenta} aspectos como a compreensão da diversidade e realidade dos \textbf{sujeitos}, das formas de produção e de trabalho e das culturas, a sustentabilidade ambiental, a diversificação da oferta de forma a possibilitar múltiplas trajetórias por parte dos estudantes e a articulação dos saberes com os contextos histórico, econômico, social, científico, ambiental, cultural local e do mundo do trabalho, especialmente articulados no artigo 5º das DCN para o ensino médio.
Tais aspectos reforçam a importância do aprofundamento oferecido pelo itinerário.
Além dos componentes \textbf{que} compõem o enfoque, outras possibilidades de articulação entre os componentes podem contribuir para o enriquecimento da formação do estudante.
Entre elas, destacam -se, a seguir : O aprofundamento da área de Matemática e suas Tecnologias contribuirá para a introdução à construção de conhecimentos técnicos necessários para as carreiras afins desse itinerário formativo, principalmente com o uso da estatística, matemática financeira e probabilidade, favorecendo \textbf{um} estudo interdisciplinar envolvendo as dimensões cultural, social, política e psicológica, além da econômica, sobre as questões de consumo, trabalho e dinheiro ( BNCC ).
Na área de Ciências da Natureza e suas Tecnologias, a ciência física dialoga com o itinerário formativo das ciências sociais e econômicas nas discussões \textbf{que} visam \textbf{contextualização} do ponto de vista \textbf{sócio-histórico}, como a produção do conhecimento científico e a evolução dos processos industriais \textbf{que} dinamizaram a vida do \textbf{homem}.
O componente curricular química interage com a área de ciências sociais, econômicas e administrativa quando aborda temáticas como consumo consciente, poluição do ar, créditos de carbono, acordos de exploração dos recursos naturais.
Ou seja, a utilização de produtos e serviços \textbf{que} ocasionam impactos políticos, econômicos, sociais, ambientais e culturais.
Neste itinerário, a biologia colabora com \textbf{tópicos} \textbf{que} auxiliem o estudante a interpretar realidades voltadas para o social, econômico e administrativo, ou seja, por meio de temas agregadores de valor econômico, seja na escala local, regional ou global, \textbf{que} perpasse por interesses sociais, como o impacto das doenças infecciosas no cenário \textbf{atual}, passando pelo estudo dos organismos \textbf{que} têm valor de mercado, assim como por suas tecnologias, culminando nas reflexões ligadas à sustentabilidade.
Percebe -se, \textbf{aqui}, \textbf{que} a compreensão do \textbf{homem} como administrador da biosfera proporciona atenção não apenas ao ser humano, mas a todas as relações existentes.
O componente curricular língua portuguesa buscará desenvolver nos cursos relacionados ao itinerário de Ciências Sociais, Econômicas e Administrativas eventos discursivos/letramentos \textbf{que} são requeridos pelos componentes curriculares do itinerário em questão, assim, buscando \textbf{aprimorar} competências e habilidades daquilo \textbf{que} o itinerário formativo busca oferecer.
O estudo da língua inglesa na \textbf{contemporaneidade}, em \textbf{um} mundo de grandes avanços tecnológicos, competitivo e de múltiplas identidades, \textbf{exige} de nossos estudantes \textbf{uma} postura cada vez mais pertencente a \textbf{um} mundo sem barreiras linguísticas e \textbf{que} lhes permita estar inseridos em \textbf{um} contexto social, tecnológico, político e econômico de cidadãos voltados para \textbf{um} mundo do mercado de trabalho mais consciente, autônomo e solidário.
Dessa forma, a importância do componente curricular de língua inglesa para o itinerário formativo de Ciências Sociais, Econômicas e Administrativas é a de contribuir para \textbf{uma} formação completa e \textbf{que} realmente faça sentido na vida prática do estudante, tanto na escola quanto na vida social e profissional \textbf{que} ele venha a escolher, aliado ao \textbf{que} o estudante traçou para seu projeto de vida.
O componente curricular educação física contribuirá, entre outras coisas, para criar itinerários formativos inovadores e flexíveis, calcados na realidade e nas necessidades locais do estudante, buscando a inclusão de todos para vivenciarem e continuarem desenvolvendo suas competências e habilidades através das práticas corporais.
Quanto à contribuição da arte no itinerário formativo para o campo das Ciências Sociais, Econômicas e Administrativas, é relevante a tomada de consciência de \textbf{que} este componente curricular possui a capacidade de potencializar a compreensão dos objetos do conhecimento, à medida \textbf{que} seu estudo \textbf{desperta} para as necessidades humanas de produção e reprodução da realidade \textbf{que} nos cerca, desde os tempos mais remotos.
Esta interação com o meio se reflete nos demais campos, onde o ser humano também buscou relacionar -se entre si, quando a arte já se faz presente nas relações humanas, sociais e econômicas desde os nossos antepassados mais remotos, nas trocas simbólicas de objetos artísticos e com valores subjetivos, racionalizados já nos estudos de Bronisław Malinowski e ainda em Pierre Bourdieu.
Reiteramos, assim, \textbf{que} a arte está presente nas áreas das ciências sociais, econômicas e administrativas – na perspectiva de potencializar a formação do \textbf{educando} nos diversos cursos de abrangência deste itinerário – \textbf{que} se dá por meio da estesia, da catarse e da fruição em arte, e, principalmente, em entender as essências e propósitos de sua inserção em cada área.
Reforçamos também \textbf{que} o processo ensino aprendizagem em arte, conforme orienta a proposta triangular de Ana Mae Barbosa, consiste em três pilares : o Fazer, o Apreciar e o Contextualizar.
Para a contribuição da área de Ciências Humanas e Sociais Aplicadas, é importante considerarmos, primeiramente, \textbf{que}, durante toda a história, a filosofia sempre foi ligada à leitura do quadro social, por meio de suas reflexões sobre a vida na pólis e seus \textbf{atores} \textbf{constituintes}.
Sendo assim, falar de filosofia sempre foi também falar de sociedade.
Diversos filósofos se ocuparam de tal empreendimento durante a história do pensamento, o \textbf{que} destaca a filosofia como a \textbf{disciplina} \textbf{que} mais contribuiu para a evolução, o desenvolvimento e a compreensão dos fenômenos sociais \textbf{inerentes} à vida humana.
A organização social, seus conceitos e problemas estão sempre ligados à construção do pensamento.
A vida na pólis antiga, em Platão e Aristóteles, por exemplo, a Cidade de Deus e a Cidade terrena, em Santo Agostinho, a Nova Sociedade, organizada pela noção de política \textbf{moderna} em Maquiavel, o biopoder, em Foucault, entre outros conceitos e \textbf{autores}, são suficientes para demonstrar a \textbf{íntima} ligação entre o filosofar e a vida em sociedade.
Sendo assim, é \textbf{imprescindível} \textbf{que} a filosofia esteja, de forma \textbf{imperativa} e preponderante, no \textbf{bojo} \textbf{central} deste itinerário, pois sua presença traz aos estudos sobre sociedade maior credibilidade, densidade e até mesmo profundidade.
Nesta perspectiva, \textbf{um} dos conceitos mais analisados pela filosofia na vida em sociedade foi, sem dúvida nenhuma, o de trabalho.
Temos grandes \textbf{autores} nesta grande heurística, como Marx, entre outros filósofos consagrados.
Nesta análise do tema trabalho, temos o aprofundamento das reflexões sobre o modo de produção e os modelos produtivos, elementos \textbf{que} são condições básicas da estrutura econômica e administrativa dentro da organização política das sociedades e da própria vida coletiva em nosso planeta.
Nesta perspectiva de \textbf{teorias} sociais, econômicas e administrativas, alheias às contribuições da reflexão filosófica, elas se tornam míopes, incompletas e/ou apenas parciais, sem \textbf{uma} visão mais ampla e totalitária dos fenômenos \textbf{que} se propõem a estudar.
Dessa forma, torna -se importante a presença da filosofia nesses estudos para \textbf{aprimorar} o processo de elaboração, descrição e construção do conhecimento pertinente a essa área de estudo.
A geografia possibilita \textbf{uma} visão \textbf{holística} sobre os processos socioespaciais, permitindo \textbf{que} o estudante desenvolva habilidades voltadas para a gestão e o planejamento territorial e econômico, com a consolidação de práticas voltadas ao desenvolvimento sustentável.
O reconhecimento das potencialidades e dos recursos territoriais auxilia na tomada de decisões de políticas públicas e possibilita a exploração de empreendimentos privados.
Nesse sentido, o estudante desenvolverá habilidades voltadas ao reconhecimento das potencialidades físicas do espaço geográfico e as relações \textbf{que} são estabelecidas com suas populações.
Essa visão ampla permitirá, por exemplo, \textbf{que} sejam estabelecidas ações de intervenção territorial, com a alocação de empreendimentos e o levantamento dos impactos a eles associados, além de permitir o gerenciamento de riscos econômicos e ambientais.
É importante, ainda, \textbf{que} se fortaleça o entendimento sobre a formação e a organização político-administrativa do Estado brasileiro.
Este conhecimento fortalecerá a compreensão da estrutura legislativa territorial e ambiental.
As análises em geografia econômica e política fortalecerão a compreensão sobre os diferentes sistemas e \textbf{teorias} econômicas, além de possibilitar a análise crítica sobre a \textbf{divisão} internacional do trabalho e a participação do Brasil no mercado internacional e sua posição na geopolítica global.
No itinerário formativo de Ciências Sociais, Econômicas e Administrativas, a ciência sociologia está presente nas discussões a respeito do mundo do trabalho e da política, analisando criticamente o desenvolvimento das estruturas sociais, principalmente neste campo produtivo, e seus impactos nas relações sociais, nas possibilidades e desafios de o Estado brasileiro gerir e administrar o país com políticas públicas eficazes, voltadas ao acesso à educação, à geração de empregos e renda, como caminhos para o desenvolvimento social coletivo e o combate às desigualdades sociais.
No eixo temático mundo do trabalho, elenca -se a \textbf{categoria} trabalho como elemento fundante da ordem social sobre a qual se desenvolveram as reflexões sociológicas dos \textbf{autores} clássicos, assim como os princípios de racionalização dessa prática.
Nesse contexto, abarcam -se \textbf{teorias} e conceitos de diferentes áreas de conhecimento em \textbf{que} o trabalho é abordado, como administração, direito, economia, história, geografia e filosofia.
No eixo temático política, compreendem -se os princípios caros à constituição do estado \textbf{moderno} e da própria administração pública, \textbf{salientando} -se questões relativas ao âmbito da política e legislação brasileiras, nas quais se articulam \textbf{teorias} e conceitos oriundos da ciência política, geografia, história, filosofia e direito.
Contextualiza -se a sociologia por meio de temáticas relativas aos seguintes objetos de conhecimento : - \textbf{Teorias} sociológicas de Durkheim, Marx e Weber ; - Princípios da administração científica : a racionalização do trabalho no \textbf{século} XX ; - A concepção de pobreza no Brasil e a distribuição da riqueza ; - A importância dos conceitos de desenvolvimento e subdesenvolvimento para análise do capitalismo no mundo \textbf{contemporâneo}.
A história contribui sobremaneira na \textbf{contextualização} do aparecimento das correntes e dos sistemas econômicos e de como tais mecanismos podem \textbf{influenciar} governos, consumo, criar \textbf{crises} e desenvolver \textbf{uma} sociedade. \textbf{Um} dos grandes exemplos está na identificação de \textbf{termos} como desenvolvimento x subdesenvolvimento, capitalismo x socialismo.
Além disso, seguir a trajetória do surgimento de correntes econômicas importantes esclarece consequências sociais, culturais e políticas de \textbf{uma} sociedade.
Com relação ao Maranhão, podemos analisar a economia de subsistência das populações quilombolas, economia solidária e a reconstrução das relações sociais de consumo, produção e trocas a partir da noção de economia e gestão.
Outra perspectiva \textbf{interessante} pode ser o desenvolvimento sustentável e solidário com enfoque territorial analisado nos aspectos de \textbf{contextualização} histórica, conceito de território rural e território de cidadania, além de experiências de promoção de desenvolvimento e uso racional dos recursos naturais nas comunidades quilombolas e comunidades de quebradeiras de coco.
Podemos, por meio da ciência econômica, analisar os períodos de recessão e crescimento econômico, a exemplo das consequências da \textbf{crise} na produção de café, o chamado “ milagre econômico ” na ditadura militar, a \textbf{crise} do petróleo no final da década de 1970, a instabilidade econômica ocorrida na Nova República e as novas perspectivas para a \textbf{humanidade} \textbf{que} envolvem o crescimento sustentável.

% matched lemmas: abordagem, aprimorar, aqui, ator, atual, autor, bojo, categoria, central, constituinte, contemporaneidade, contemporâneo, contextualização, crise, despertar, disciplina, divisão, educar, enquanto, exigir, fomentar, holístico, homem, humanidade, imperativo, imprescindível, inerente, influenciar, interessante, moderno, que, salientar, sujeito, século, sócio-histórico, teoria, termos, tópico, um, íntimo
\end{textsample}
